\documentclass[11pt,a4paper]{article}
\usepackage[margin=1in]{geometry}
\usepackage{booktabs,hyperref,graphicx,amsmath,amssymb}
\usepackage[numbers]{natbib}

\newcommand{\synth}{\textsc{synthetic}}

\title{How Long Is a Presence?\\Refuting a Geometric Derivation of the Alert\\Collapse Window in Multi-Analytics CCTV}
\author{Amit Dua\\Yushu Excellence Technologies Pvt.\ Ltd.\\\texttt{mail.amitduabits@gmail.com}}
\date{2026}

\begin{document}
\maketitle

\begin{abstract}
A CCTV platform running several analytics over the same scene emits one alert per
detection unless it collapses them. The collapse predicate --- same entity, same
window --- is standard: it is \citet{debar2001aggregation}'s aggregation predicate
under a normalised record schema of the IDMEF \citep{idmef2007} and OCSF
\citep{ocsf} lineage, evaluated over keyed session windows in the sense of
\citet{akidau2015dataflow}. The predicate is not our contribution and we do not
claim it. The open parameter is the window $W$, and our deployment set it to
$120$\,s with a derivation: camera field-of-view depth divided by permitted speed.

\textbf{That derivation is wrong, and our own data refutes it.} Computed per
camera across a 200-camera estate, the geometric dwell time is $1.27$ to
$9.63$\,s with a median of $3.66$\,s --- one and a half orders of magnitude below
the deployed constant. The deployed $120$\,s over-collapses relative to its own
stated justification at every camera in the estate ($200$ of $200$). Worse, the
geometric quantity is not the right one: what the window must span is the interval
between successive \emph{observations} of an entity, which includes revisits and
re-detections, not the time to traverse one field of view.

Measured on the operating curve, the knee sits at $15$--$30$\,s. Moving from
$15$\,s to $120$\,s buys a further $2.8\%$ reduction in alerts and costs $337$
additional masked incidents, taking distinct-incident recall from $0.9961$ to
$0.9680$. We give the corrected procedure: choose $W$ from the
inter-observation-time distribution by a cost-weighted operating point, with
\citet{satopaa2011kneedle} as the parameter-free default. Separately, we confirm
that an entity-keyed predicate spanning modalities beats per-modality
deduplication by $10.5\%$ to $43.5\%$ as the dual-tagged incident share rises from
$0.15$ to $1.0$, at no recall cost, and that a confidence-voting alternative
achieves the largest raw reduction ($88.1\%$) only by dropping $63.5\%$ of
incidents and is not a usable operating point.

All results are \synth{}, from a seeded generator over $59{,}984$ detections and
$12{,}000$ incidents. \S\ref{sec:limits} states what real operator-marked ground
truth must show.
\end{abstract}

\section{Introduction}

Run automatic number-plate recognition, face matching and region-occupancy
detection over the same camera and the same vehicle generates three detections a
second. Emit one alert each and the control room receives five alerts per
incident; the literature on alert fatigue \citep{tariq2025alertfatigue,
alahmadi2022falsepositives, sundaramurthy2015burnout} and the older process-control
work on alarm floods \citep{wang2016alarmoverview, isa182, eemua191} is unanimous
that this is the failure mode that ends the deployment.

The fix is to collapse: emit one alert per entity per window. The question that
decides whether the fix works is how long the window is. Too short and the
operator gets duplicates; too long and a genuinely new incident involving the same
entity is masked by a still-open alert. This is a two-sided error and it has an
operating curve.

Our deployment chose $W = 120$\,s and documented a derivation for it: a camera's
field of view is a few tens of metres deep, permitted speed is a few tens of
km/h, so a vehicle is present for on the order of a minute or two. This paper
reports that the derivation does not survive contact with the estate's own
parameters, explains why the geometric quantity is the wrong one in principle as
well as in magnitude, and replaces it with a procedure.

\subsection{What is and is not claimed}

Not claimed: the normalised alert record (IDMEF \citep{idmef2007}, OCSF
\citep{ocsf}, JDL level-2 common referencing \citep{steinberg1999jdl,
hall1997introduction}); the same-entity-same-window aggregation predicate
\citep{debar2001aggregation, julisch2003clustering, valdes2001probabilistic};
keyed session windowing \citep{akidau2015dataflow, cugola2012processing}; or
duplicate detection as a discipline \citep{fellegi1969theory,
elmagarmid2007duplicate}. All of this is background and \S\ref{sec:bg} says so in
its first sentence.

Claimed: (i) the geometric derivation of $W$ is refuted by the estate's own
geometry, quantitatively (\S\ref{sec:refute}); (ii) the measured knee is
$15$--$30$\,s and the deployed value sits far past it, with the cost quantified
(\S\ref{sec:eval}); (iii) a defensible replacement procedure for choosing $W$
(\S\ref{sec:choose}); (iv) the cross-modal entity-keyed predicate's marginal value
over per-modality deduplication, measured against the dual-tagging rate
(\S\ref{sec:crossmodal}).

\section{Background, all of it established}
\label{sec:bg}

\emph{Nothing in this section is a contribution.} Alert normalisation into a
common record with entity, time, source and confidence fields is IDMEF
\citep{idmef2007} and, in current practice, OCSF \citep{ocsf}; the fusion framing
is JDL level 2 \citep{steinberg1999jdl}, surveyed by
\citet{khaleghi2013multisensor}. Aggregating alerts that share a key within a
time window is \citet{debar2001aggregation}, refined by
\citet{julisch2003clustering} and \citet{ning2002constructing}, with a
probabilistic variant in \citet{valdes2001probabilistic}. Evaluating such a
predicate over unbounded streams with per-key session windows is
\citet{akidau2015dataflow}. Session-gap estimation from inter-event times has its
own literature in web analytics \citep{catledge1995characterizing,
halfaker2015session, meiss2009session}, and in process control the same problem is
the alarm dead-band and delay-timer design question
\citep{adnan2011detectiondelay, afzal2018timedeadbands, wang2022delaytimers}.

We adopt all of it. The gap is that none of these sources tells a CCTV operator
what number to type in, and the one derivation our deployment invented for itself
is wrong.

\section{The collapse predicate as deployed}

A detection is a tuple $(e, m, c, t, s)$: entity key, modality, camera, timestamp,
score. The predicate is entity-keyed and modality-agnostic:

\begin{quote}
Emit an alert for detection $d$ iff no alert with the same entity key $e$ is open;
an alert opened at $t_0$ is open over $[t_0, t_0 + W)$.
\end{quote}

Window semantics are \emph{anchored}, not merging: a detection inside an open
window does not extend it. This matters and is under-reported in the literature.
Merging semantics turn a stream of detections at a busy junction into one
unbounded alert; anchored semantics bound the suppression at $W$ by construction.
Every number in this paper is under anchored semantics.

Two error modes follow. A \emph{redundant alert} is a second alert for an incident
already alerted. A \emph{masked incident} is a distinct incident that produced no
alert because an unrelated earlier incident with the same entity key held the
window open. We report distinct-incident recall for the second, because plain
incident recall is $1.0$ everywhere and hides the failure.

\section{Why the geometric derivation fails}
\label{sec:refute}

\begin{table}[t]
\centering\small
\begin{tabular}{lr}
\toprule
\textbf{Quantity} & \textbf{Value} \\
\midrule
Geometric dwell, minimum over estate & 1.27\,s \\
Geometric dwell, median over estate & 3.66\,s \\
Geometric dwell, maximum over estate & 9.63\,s \\
Deployed constant $W$ & 120\,s \\
Cameras where the constant exceeds the geometric value & 200 of 200 \\
\bottomrule
\end{tabular}
\caption{\textbf{Geometric window per camera against the deployed constant
(\synth, 200 cameras).} $W_{\text{geom}} = \text{FOV depth} / \text{permitted
speed}$. The deployed constant exceeds its own stated derivation by a factor of
$12$ at the most generous camera and $94$ at the tightest.}
\label{tab:geom}
\end{table}

Table~\ref{tab:geom} is the refutation. Field-of-view depth divided by permitted
speed gives seconds, not minutes, at every camera in the estate. If the deployment
had applied its own stated rule it would have set $W$ between $1$ and $10$\,s, and
per camera rather than globally. It did not; it set $120$\,s. The derivation was
therefore not the reason for the value, whatever the design document says.

The magnitude error is the smaller problem. The conceptual error is that dwell
time is not the quantity the window must span. The window exists to suppress
\emph{repeat observations of the same entity that belong to the same operational
event}. Those arise from (a) multiple analytics firing on one pass, (b) repeated
detections across frames during one pass, and (c) the entity leaving and returning
--- circling a block, stopping and restarting, being observed by an adjacent
camera whose alerts share the entity key. Case (c) dominates the tail and has
nothing to do with how long one field of view is. A derivation from geometry
cannot see it. The correct input is the empirical inter-observation-time
distribution.

\section{Choosing $W$ properly}
\label{sec:choose}

Four families of method apply, and the choice among them is a value judgement that
should be made explicitly rather than buried in a constant.

\textbf{Cost-weighted operating point.} Treat window selection as threshold
selection on a two-sided error curve \citep{fawcett2006roc}. Assign a cost
$\lambda_r$ to a redundant alert (operator seconds) and $\lambda_m$ to a masked
incident (a missed event), and minimise $\lambda_r R(W) + \lambda_m M(W)$. This is
the framing we lead with, because it makes the deployment's values legible and
lets another deployment substitute its own. \citet{axelsson2000baserate} is the
reason the two costs are not comparable by intuition.

\textbf{Mixture crossover.} Fit a two-component mixture to inter-observation times
--- within-event and between-event --- and set $W$ at the posterior crossover. This
is the web-analytics session-gap method \citep{catledge1995characterizing,
meiss2009session} and is the most principled where the two populations are
separable.

\textbf{Knee detection.} Apply \citet{satopaa2011kneedle} to the
reduction-versus-masking curve. Parameter-free, and the right default where no
cost model exists.

\textbf{Spec-based timer design.} The process-control approach
\citep{adnan2011detectiondelay, afzal2018timedeadbands, wang2022delaytimers}: set
the dead-band from a specified maximum acceptable detection delay. Appropriate
where a regulator states the delay bound.

All four are defensible. Dividing a distance by a speed is not one of them.

\section{Evaluation}
\label{sec:eval}

\begin{table}[t]
\centering\small
\begin{tabular}{rrrrr}
\toprule
\textbf{$W$ (s)} & \textbf{Alerts} & \textbf{Reduction vs.\ $W{=}0$} & \textbf{Masked incidents} & \textbf{Distinct-incident recall} \\
\midrule
0    & 59{,}984 & ---     & 0     & 1.0000 \\
5    & 12{,}962 & 78.4\% & 19    & 0.9984 \\
15   & 11{,}961 & 80.1\% & 47    & 0.9961 \\
30   & 11{,}915 & 80.1\% & 87    & 0.9928 \\
60   & 11{,}844 & 80.3\% & 164   & 0.9863 \\
120  & 11{,}626 & 80.6\% & 384   & 0.9680 \\
180  & 11{,}345 & 81.1\% & 672   & 0.9440 \\
240  & 11{,}105 & 81.5\% & 904   & 0.9247 \\
480  & 10{,}378 & 82.7\% & 1{,}630 & 0.8642 \\
900  & 9{,}288  & 84.5\% & 2{,}716 & 0.7737 \\
1800 & 9{,}012  & 85.0\% & 2{,}988 & 0.7510 \\
\bottomrule
\end{tabular}
\caption{\textbf{The operating curve (\synth, $59{,}984$ detections,
$12{,}000$ incidents).} Redundant alerts are already down to $8$ at $W{=}15$\,s
from $47{,}984$ at $W{=}0$. Everything bought beyond $15$\,s is bought from
recall.}
\label{tab:sweep}
\end{table}

Table~\ref{tab:sweep} is the paper's central result. Read it in two parts.

\textbf{Below the knee, collapse is nearly free.} Going from $0$ to $15$\,s
removes $47{,}976$ of the $47{,}984$ redundant alerts and masks $47$ incidents:
$80.1\%$ fewer alerts for $0.4\%$ of distinct-incident recall. No deployment
should decline this.

\textbf{Above the knee, the curve is almost flat in the benefit and steep in the
cost.} From $15$\,s to $120$\,s, alerts fall by a further $335$ ($2.8\%$ of the
remainder, $0.5$ percentage points of total reduction) while masked incidents rise
from $47$ to $384$, an eightfold increase. Recall drops from $0.9961$ to $0.9680$.
The deployed value therefore sits well past the knee: it pays $337$ masked
incidents for a reduction an operator would not notice. At $1{,}800$\,s the trade
is grotesque --- $85.0\%$ reduction against $0.7510$ recall --- but the shape of the
curve is already established by $120$\,s.

\citet{satopaa2011kneedle} applied to the (reduction, masking) curve returns a knee
in $15$--$30$\,s, which is where we recommend the deployment move, subject to the
cost weighting in \S\ref{sec:choose}.

\subsection{Cross-modal collapse against per-modality deduplication}
\label{sec:crossmodal}

\begin{table}[t]
\centering\small
\begin{tabular}{lrrrrr}
\toprule
\textbf{Method} & \textbf{Alerts} & \textbf{Reduction} & \textbf{Incident recall} & \textbf{Distinct recall} & \textbf{Alerts/incident} \\
\midrule
naive (one per detection) & 59{,}984 & ---     & 1.000 & 1.000 & 5.00 \\
per-modality dedup        & 15{,}788 & 73.7\% & 1.000 & 0.9635 & 1.32 \\
confidence voting         & 7{,}130  & 88.1\% & \textbf{0.365} & 0.3483 & 1.63 \\
entity-agnostic collapse  & 11{,}626 & 80.6\% & 1.000 & 0.9680 & 0.97 \\
\bottomrule
\end{tabular}
\caption{\textbf{Collapse strategies (\synth), $W = 120$\,s.} Confidence voting
achieves the largest reduction by discarding $63.5\%$ of incidents entirely. It is
reported because it is the obvious alternative, and rejected because no control
room can run at $0.365$ incident recall.}
\label{tab:baselines}
\end{table}

\begin{table}[t]
\centering\small
\begin{tabular}{rrrr}
\toprule
\textbf{Dual-tagged share} & \textbf{Per-modality alerts} & \textbf{Entity-agnostic alerts} & \textbf{Extra reduction} \\
\midrule
0.00 & 11{,}654 & 11{,}654 & 0.0\% \\
0.15 & 12{,}990 & 11{,}633 & 10.5\% \\
0.30 & 14{,}415 & 11{,}671 & 19.0\% \\
0.45 & 15{,}715 & 11{,}605 & 26.2\% \\
0.70 & 17{,}973 & 11{,}637 & 35.3\% \\
1.00 & 20{,}614 & 11{,}644 & 43.5\% \\
\bottomrule
\end{tabular}
\caption{\textbf{Marginal value of crossing modalities (\synth).} Recall is
$1.000$ for both methods at every share. The entity-agnostic column is flat
because it is insensitive to how many modalities tagged the incident; that
insensitivity is the entire benefit.}
\label{tab:crossmodal}
\end{table}

Table~\ref{tab:crossmodal} isolates the one design decision in the predicate that
is genuinely ours to defend: keying on the entity rather than on
(entity, modality). The benefit is exactly proportional to how often an incident is
seen by more than one analytic, from zero at $0\%$ dual tagging to $43.5\%$ at
$100\%$. Our estate runs at roughly $40\%$ dual tagging
($4{,}859$ of $12{,}000$ incidents), placing the real benefit near $24\%$. A
deployment running one analytic per camera should not adopt this and gains
nothing from it; that conditional is the honest form of the claim.

\section{Related work}

Covered in \S\ref{sec:bg} and conceded there. The one comparison we owe is
\citet{valdes2001probabilistic}, whose weighted-similarity aggregation is the
principal alternative to a hard entity key: it fuses on graded feature agreement
rather than exact match. It would handle entity-key errors --- a misread plate ---
which our predicate cannot, and a deployment with noisy keys should prefer it. We
do not claim superiority over it; our estate's keys come from ANPR with a
confidence gate, and exact matching is adequate there.
\citet{cheng2013smithwaterman} and \citet{ye2022reid} are the alternatives for
graded matching across cameras.

\section{Limitations}
\label{sec:limits}

\begin{enumerate}
\item \textbf{Synthetic ground truth.} ``Distinct incident'' is a label our
generator emits. In an estate, whether two appearances of one vehicle forty
seconds apart are one incident or two is an operator judgement, and the whole
curve in Table~\ref{tab:sweep} depends on that judgement. The replication that
matters is operator-marked incident boundaries on real footage; if operators mark
a longer typical event than our generator, the knee moves right and $120$\,s may
be defensible on grounds we did not test. The geometric refutation in
\S\ref{sec:refute} does not depend on this, since it compares the deployment
against its own stated rule.
\item \textbf{Anchored semantics.} Under merging semantics all of these numbers
change and the masking cost grows without bound.
\item \textbf{One estate, one entity-key source.} Key errors are not modelled.
\item \textbf{A single global $W$.} Table~\ref{tab:geom} shows the geometric
quantity varies $7.6\times$ across cameras. Whether a per-camera $W$ pays is not
answered here.
\item \textbf{Case grouping is a different problem.} An operator may want all
alerts for one entity grouped in the interface for hours while still being
notified of each. That is a presentation decision and it is arguably the right
answer to the fatigue problem; this paper addresses suppression, not grouping.
\end{enumerate}

\section{Conclusion}

The collapse predicate is not new and we do not claim it. Its window is the
parameter that decides whether it helps, and the derivation our deployment used
for that window --- field-of-view depth over permitted speed --- is refuted by the
estate's own camera parameters at all $200$ cameras, by a factor of $12$ to $94$,
and is measuring the wrong quantity besides. On the operating curve the knee is at
$15$--$30$\,s; the deployed $120$\,s pays $337$ additional masked incidents for a
$2.8\%$ further reduction in alerts. We recommend deployments derive $W$ from the
inter-observation-time distribution under an explicit cost weighting, publish the
weighting, and stop deriving it from geometry.

\section*{Reproduction}

\begin{verbatim}
cd code
pip install -r requirements.txt
python -m pytest -q
python -m prresearch.p5_fusion.experiment
\end{verbatim}

Fixed seeds; \texttt{results/p5\_fusion.json} is byte-identical across runs.
Code and data: \url{https://github.com/amitduabits/PRAHARI-P5-EventSchema}

\bibliographystyle{plainnat}
\bibliography{references}

\end{document}
